\chapter{Microestructura y ejecución}\label{cap:micro}

\section{El libro de órdenes}
Un mercado electrónico es un \textbf{libro de órdenes limitadas}: a cada precio hay una cola de
órdenes de compra (\emph{bid}) y de venta (\emph{ask} u \emph{offer}). El mejor bid y el mejor ask
forman la horquilla (\emph{spread}); el precio medio es $m = (b+a)/2$. Quien envía una orden a mercado
\emph{cruza} el spread y paga liquidez; quien deja una orden limitada la \emph{provee} y espera.

\begin{table}[h]\centering\small
\begin{tabularx}{\textwidth}{lX}
\toprule
Orden & Uso \\ \midrule
Mercado & Ejecución inmediata al precio disponible. Riesgo: deslizamiento (\emph{slippage}) en libros finos o en noticias. \\
Limitada & Precio máximo (compra) o mínimo (venta). Riesgo: no ejecutarse justo cuando el precio se va. \\
Stop (stop-market) & Se convierte en orden a mercado al tocar el nivel. Garantiza salida, no precio (\emph{gaps}). \\
Stop-limit & Se convierte en limitada al tocar el nivel. Garantiza precio, no salida: peligrosa para stops de protección. \\
OCO / bracket & Objetivo y stop enlazados: si se ejecuta uno se cancela el otro. \\
Iceberg & Muestra solo una parte del tamaño. \\
Validez & DAY, GTC (hasta cancelar), IOC (lo que se pueda ya), FOK (todo o nada), MOC/LOC (al cierre). \\
\bottomrule
\end{tabularx}
\caption{Tipos de orden habituales.}
\end{table}

\section{Costes de transacción}
El coste total de una operación tiene tres partes:
\begin{equation}
\text{coste} = \underbrace{\tfrac12\,\text{spread}}_{\text{explícito}} + \text{comisiones} + \underbrace{\text{impacto}}_{\text{implícito}} .
\end{equation}
El impacto de mercado de una orden grande crece aproximadamente con la \textbf{raíz cuadrada} del
tamaño relativo al volumen diario (ley empírica de la raíz cuadrada):
\begin{equation}\label{eq:impacto}
\frac{\Delta P}{P} \approx Y\,\sigma_{\text{diaria}}\sqrt{\frac{Q}{V}},\qquad Y\approx 0{,}5\text{--}1,
\end{equation}
donde $Q$ es la cantidad y $V$ el volumen diario. Ejemplo: con $\sigma_{\text{diaria}}=2~\%$ y una orden
del 1~\% del volumen diario, el impacto ronda $0{,}1$--$0{,}2~\%$: \emph{más que el spread} de una acción
líquida. Por eso las órdenes grandes se trocean con algoritmos:
\begin{itemize}
\item \textbf{TWAP}: reparte la orden uniformemente en el tiempo.
\item \textbf{VWAP}: reparte según el perfil de volumen intradía (más al principio y al final de la sesión).
\item \textbf{POV}: participa en un porcentaje fijo del volumen que va apareciendo.
\item \textbf{Implementation shortfall}: minimiza la diferencia con el precio de decisión, equilibrando impacto y riesgo de esperar.
\end{itemize}

\begin{practica}
Mide tu ejecución: anota el \emph{precio de decisión} (cuando decides operar) y el precio medio de
ejecución. La diferencia, en puntos básicos, es tu \emph{slippage}. En el backtest
(\py{trader.backtest.backtest}, argumento \py{cost\_bps}) usa un coste igual o mayor al que realmente
pagas; si la estrategia deja de funcionar con el doble de coste, no tiene margen de seguridad.
\end{practica}

\section{Sesiones y liquidez}
\begin{itemize}
\item \textbf{Acciones}: subasta de apertura, sesión continua, subasta de cierre (el cierre concentra
  el 10--20~\% del volumen en EE.~UU. y es el precio de referencia de índices y fondos). Pre y
  post-mercado con poca liquidez.
\item \textbf{FX}: mercado OTC 24 horas de domingo noche a viernes noche; liquidez máxima en el solape
  Londres--Nueva York (aprox. 12--16~h UTC). Fixing WM/Reuters de las 16:00 de Londres.
  \py{trader.fx.open\_sessions(hora\_utc)} dice qué sesiones están abiertas.
\item \textbf{Futuros}: casi 24 horas en Globex/ICE, pero el volumen se concentra en el horario del
  mercado físico (\emph{pit hours}). El precio de liquidación diario (\emph{settlement}) fija el margen.
\item \textbf{Cripto}: 24/7, con liquidez fragmentada entre plataformas.
\end{itemize}

\section{Apalancamiento y margen}
En futuros, CFD, FX al contado con broker y perpetuos de cripto no pagas el nocional, sino un
\textbf{margen inicial} (típicamente 3--15~\% del nocional). Si las pérdidas reducen la cuenta por debajo
del \textbf{margen de mantenimiento}, el broker pide más fondos (\emph{margin call}) o liquida. El
apalancamiento efectivo de la cuenta es
\[
\text{apalancamiento} = \frac{\sum_i |\text{nocional}_i|}{\text{capital}} ,
\]
y es lo que hay que vigilar, no el margen exigido: un margen del 5~\% permite 20x, lo que no significa
que debas usarlo. \py{scripts/informe\_riesgo.py} lo calcula para tu cartera.

\begin{ejercicios}
\begin{enumerate}
\item Una acción cotiza a 49,98 / 50,02. Compras 2000 acciones a mercado y vendes a mercado al cabo de
  un rato sin que el precio medio se mueva. ¿Cuánto has perdido por el spread? ¿Y en pb del nocional?
\item Con la ecuación~\eqref{eq:impacto} ($Y=0{,}7$, $\sigma=1{,}8~\%$) estima el impacto de comprar el
  5~\% del volumen diario. ¿Cuántos días tendrías que repartir la orden para que el impacto diario sea
  menor que 10 pb?
\item ¿Por qué un stop-limit puede no protegerte en un gap de apertura?
\end{enumerate}
\end{ejercicios}
